\section{Optimización compleja y codificación de restricciones en DSPy}

La última parte describe cómo el marco DSPy maneja las restricciones no lineales difíciles de expresar en la optimización combinatoria, como el diseño birefringent en nanophotonics.

\subsection{Linealización latente \& Green Descent}

Para las restricciones complejas, el equipo primero predice un cost vector \(C\), luego convierte el problema original en un subproblema lineal y, finalmente, usa un solver para generar \(X^\star\). Todo el proceso forma un grafo de cómputo que va de la description \(Y\) a \(C(Y)\) y luego a la salida del solver; el loss es la diferencia respecto al objetivo original en el dominio temporal, y se retropropaga mediante gradient descent para entrenar la representación de \(C\).
\begin{figure}[H]
\centering
\includegraphics[width=\textwidth]{frame-latent-opt.jpg}
\caption{Grafo de cómputo en DSPy que deriva \(C\) a partir de la descripción \(Y\) y luego llama al solver para producir \(X^\star\). \protect\footnotemark}
\end{figure}
\footnotetext{Intervalo de tiempo del video: 00:45:05--00:45:25.}
\begin{knowledgebox}{Latent Optimization Pipeline}
1) A partir de la descripción se predice el latent cost \(C\); 2) con \(C\) fijo, se resuelve de manera aproximada con un solver diferenciable; 3) la salida se reinserta en el loss físico; 4) se entrena \(C\) con descenso de gradiente. Esta resolución inversa, veloz como un relámpago, es más estable que resolver directamente el problema no lineal original.
\end{knowledgebox}
\begin{importantbox}{Pasos jerárquicos de Green Descent}
Primero se usa un solver ligero para resolver analíticamente el subproblema linealizado y generar \(X^\star\); después la salida pasa por backprop de regreso al latent cost, y con descenso de gradiente se ajusta \(C\). Esta iteración de dos niveles (solver + actualización de gradiente) evita numéricamente las oscilaciones que produce ascender directamente por el espacio no lineal original.
\end{importantbox}

\subsection{Restricciones ópticas y de fabricación}

En el diseño óptico de una cuadrícula de 80\u00d780, cada punto de la cuadrícula solo puede tomar el valor 0/1, y el proceso de fabricación exige que no aparezcan puntos aislados y que se respete el tamaño del pincel. También hay que considerar restricciones multidimensionales como la respuesta en frecuencia y las trayectorias de interferencia, por lo que el problema se descompone y se valida en varios benchmark (beam splitter, wavelength multiplexer).
\begin{warningbox}{Trampas ocultas de las restricciones de fabricación}
Las restricciones no solo tienen forma algebraica, sino que también incluyen la fabricabilidad de ingeniería: las restricciones lineales derivadas del tamaño del pincel, los píxeles aislados inviables y las limitaciones topológicas de las trayectorias de la luz. Cualquier optimización que ignore esto resultará inservible en la etapa de fabricación.
\end{warningbox}

\subsection{Resumen de la sección}

Mediante la combinación de latent linearization y la retroescritura del solver, DSPy logra que el aprendizaje de restricciones complejas pase de un ``espacio de búsqueda infinito'' a ``parámetros diferenciables'', y a la vez incorpora las restricciones de fabricación al sistema de evaluación.

